\section{Laplacetransformation und Differentialgleichungen}\label{MA2-laplacetransformation-und-differentialgleichungen}

\begin{table}[H]
  \centering
  \begin{tabularx}{\textwidth}{@{}>{\raggedright\arraybackslash}X r@{}}
    \toprule
    \multicolumn{2}{c}{\textbf{�bersicht �ber Lernziele und Voraussetzungen}}\\
    \midrule
    \textbf{Lernziele} & \\[.5\normalbaselineskip]
    \tabitem lineare Anfangswertprobleme mit Hilfe der Laplace-Transformation l�sen & \\
    \tabitem das Transformationsverfahren systematisch anwenden (Transformieren, algebraisch l�sen, R�cktransformieren) & \\[.5\normalbaselineskip]
    \midrule
    \textbf{Voraussetzungen} &  \\[.5\normalbaselineskip]
    \tabitem die Wirkung der Laplace-Transformation auf Ableitungen als Regel anwenden k�nnen & \abschnittsverweis{MA2-laplacetransformation} \\
    \tabitem Partialbruchzerlegungen durchf�hren k�nnen & \abschnittsverweis{MA2-partialbruchzerlegung} \\[.5\normalbaselineskip]
    \bottomrule
  \end{tabularx}
\end{table}

\noindent \ifintegriereaufgaben Die �bungsaufgaben zum Thema finden Sie \hyperref[MA2-laplacetransformation-und-differentialgleichungen-aufgaben]{hier}.
\vspace{0.3cm}
\hrule
\fi

\subsection{Idee}

Die Eigenschaften der Laplace-Transformation bzgl. \underline{Ableiten} und \underline{Linearit�t}\randverweis{MA2-laplacetransformation} legen nahe die Laplace-Transfor\-ma\-tion zur L�sung von \underline{gew�hnlichen Differentialgleichungen}\randverweis{MA2-differentialgleichungen} anzuwenden.

Gegeben sei Anfangswertproblem mit einer Differentialgleichung der Form
\[
	\ddot y(t)+5 \dot y(t)+6y(t) = t^2\,.
\]
Die Differentialgleichung ist inhomogen, linear, von zweiter Ordnung und hat nur konstante Koeffizienten. Die Idee ist beide Seiten der Differentialgleichung einer Laplace-Transformation zu unterziehen. Was passiert?

\subsection{Vorgehensweise}

\begin{exabox}{Anfangswertproblem Part $I$}{\hspace{-0.4cm}$\textnormal{AWP} \to Y(s)$ }
	Anfangswertproblem:
\[
	\begin{cases}
		y(0) = y_0\,, \quad & t=0\,,\\
		\dot y(0) = y_1\,, \quad & t=0\,,\\
	\ddot y(t)+5 \dot y(t)+6y(t) = t^2\,,\quad & t>0\,.\end{cases}
\]
Laplace-Transformation:

\begin{eqnarray*}
	\begin{array}{c rcl}
		& \left( -\dot y(0) - sy(0) + s^2Y(s) \right) + 5\left( -y(0) + sY(s)\right) + 6 Y(s) &=& 2 s^{-3}  \\[2ex]
		\Leftrightarrow & \left( s^2+5s+6\right) Y(s) -\dot y(0) - (5+s)y(0) &=& 2 s^{-3} \\[2ex]
	\Leftrightarrow & Y(s) &=& \displaystyle{\frac{2s^{-3} + \dot y(0) + (5+s)y(0)}{s^2+5s+6}}\,.
\end{array}
	\end{eqnarray*}
	L�sung im Bildraum:
	\[ Y: s\mapsto \frac{2s^{-3} + \dot y(0) + (5+s)y(0)}{s^2+5s+6} \,.\]
\end{exabox}

Und damit haben wir eine Differentialgleichung durch Laplace-Transformation in eine \underline{algebraische Gleichung}\randverweis{MA1-algebraische-gleichungen} f�r die Laplace-Transformierte $Y$ der L�sung transformiert! Wir wissen nun also wie die Laplace-Transformierte unserer L�sung aussieht:
$$ Y(s) = \displaystyle{\frac{2s^{-3} + \dot y(0) + (5+s)y(0)}{s^2+5s+6}}\,. $$

Leider sind wir nicht an $s\mapsto Y(s)$, sondern an der L�sung $y\mapsto y(t)$ im Zeitbereich interessiert und deshalb stellen sich die folgenden Fragen:
\begin{itemize}
	\item K�nnen wir zur�ck transformieren?
	\item Falls ja, ist die R�cktransformation eindeutig bestimmt?
\end{itemize}

Anstatt eine inverse Laplace-Transformation zu definieren und durchzuf�hren\footnote{Bei der Fourier-Transformation ist der Weg zur�ck klar:$\int\limits_{-\infty}^\infty \hat f(\omega) e^{i\omega t}\,\mathrm{d}\omega$.}, nutzen wir die Struktur der \underline{rationalen Funktion}\randverweis{MA1-rationale-funktionen} aus! Die Idee ist, eine Tabelle mit Laplace-Transformationen r�ckw�rts zu lesen. Zuvor muss die komplizierte rationale Funktion in einfache Teile zerlegt werden, die wir in der Tabelle auch finden. Das bedeutet praktisch, dass eine \underline{Partialbruchzerlegung}\randverweis{MA2-partialbruchzerlegung} durchzuf�hren ist.\marginpar{\small Aufgaben}

Bez�glich der zweiten Frage, merken wir uns: Es gibt nur eine Funktion, deren Laplace-Transformation dieses $Y$ ist. Wenn man diese Funktion gefunden hat, hat man die L�sung der Differentialgleichung gefunden!
%Man beachte: Die Laplace-Transformation ist linear: Summen von Funktionen macht sie zu Summen ihrer Transformierten; ein konstantes Vielfaches einer Funktion macht sie zu diesem Vielfachen deren Transformierter.
\begin{exabox}{Anfangswertproblem Part $II$}{\hspace{-0.4cm}$Y(s)\to y(t)$ }
	L�sung im Bildraum:
\[
	Y(s) =\frac{2s^{-3} + \dot y(0) + (5+s)y(0)}{s^2+5s+6}
\]
L�sung im Zeitbereich:
	\begin{enumerate}
		\item $Y(s)$ wird so umgeformt, dass im Z�hler und im Nenner Polynome in $s$ stehen:
			\[
				Y(s) =\frac{2 + s^3 \dot y(0) + s^3 (5+s)y(0)}{s^3\cdot( s^2 + 5s+6 )}
			\]
		\item Der Nenner von $Y(s)$ wird weitestgehend in Faktoren zerlegt (Nullstellensuche und Polynomdivision):
			\[
				s^3\cdot (s^2 + 5s+6)= s^3\cdot (s+3)\cdot (s+2)\,.
			\]
		\item F�r $Y(s)$ wird eine Partialbruchzerlegung mit dem Ansatz
			\[
				Y(s) = \frac{A_1}{s} + \frac{B_1}{s^2}+\frac{C_1}{s^3} + \frac{A_2}{s+3} + \frac{A_3}{s+2}
			\]
			durch und erh�lt durch Koeffizientenvergleich die folgende Zerlegung
			\[
				Y(s)=\frac{1}{3 s^3}- \frac{5}{18s^2} + \frac{19}{108s} +\frac{-\dot y(0)-2y(0)+2/27}{s+3} + \frac{\dot y(0)+3y(0)-1/4}{s+2}\,.
			\]
		\item F�r jeden der f�nf Teile schl�gt man nach, welche Funktion hier jeweils transformiert wurde. Damit hat man die L�sung der Differentialgleichung gefunden:
			\[
				y(t) = \frac{1}{6}t^2 - \frac{5}{18} t + \frac{19}{108} + \left( -\dot y(0) - 2y(0) + \frac{2}{27} \right) \cdot e^{-3t} + \left( -\dot y(0) + 3y(0) - \frac{1}{4} \right) \cdot e^{-2t}\,.
			\]
	\end{enumerate}
\end{exabox}




